\documentclass{article}
\usepackage{graphicx} % Required for inserting images
\usepackage[a4paper, left=2cm, right=2cm, top=3cm, bottom=3cm]{geometry}
\usepackage{amsmath}
\usepackage{amssymb}
\usepackage{listings}
\usepackage{bm}
\usepackage{xcolor}
\lstset{
  basicstyle=\ttfamily\small,
  keywordstyle=\color{blue},
  stringstyle=\color{red},
  commentstyle=\color{gray},
  showstringspaces=false,
  frame=single
}


\title{Mathematical Foundations PS7}
\author{Siheon Lee}
\date{November 2025}

\begin{document}

\maketitle

\section{Chapter 6 Exercise 1 (pg. 110)}
\begin{flalign*}
    a \cdot a &= 1 \cdot 1 + 1 \cdot 1 = 2 \\
    \vert \vert a \vert \vert &= \sqrt{a \cdot a} = \sqrt{2}
\end{flalign*}

\section{Chapter 6 Exercise 2 (pg. 110)}
\begin{flalign*}
    x &= \left[ \begin{array}{cc}
        3 & 1 
    \end{array}\right] \\
    b &= \left[ \begin{array}{cc}
        1 & 1
    \end{array} \right] \\
    \vert \vert b \vert \vert &= \sqrt{2} \\
    \hat{b} &= \frac{b}{\vert \vert b \vert \vert} = \left[ \begin{array}{cc}
        \frac{1}{\sqrt{2}} & \frac{1}{\sqrt{2}}
    \end{array} \right] \\
    x \cdot \hat{b} &= \frac{3}{\sqrt{2}} + \frac{1}{\sqrt{2}} = \frac{4}{\sqrt{2}} = 2\sqrt{2} \\
    Proj_{b}(x) &= (x \cdot \hat{b})\hat{b} = \left[ \begin{array}{cc}
        2 & 2
    \end{array} \right]
\end{flalign*}

\section{Chapter 6 Exercise 3 (pg. 110)}
\begin{flalign*}
    \vert \vert a \vert \vert \vert \vert b \vert \vert \cos(\theta) &= \vert \vert b \vert \vert (a \cdot \hat{b}) \\
    &= \vert \vert b \vert \vert (a \cdot \frac{b}{\vert \vert b \vert \vert}) \\
    &=\frac{\vert \vert b \vert \vert}{\vert \vert b \vert \vert} (a \cdot b) \\
    &= a \cdot b = \left[ \begin{array}{cc}
        a_1 & a_2
    \end{array} \right] \cdot \left[ \begin{array}{cc}
        b_1 & b_2
    \end{array} \right] = \sum_{i=1}^2 a_ib_i
\end{flalign*}

\section{Chapter 6 Exercise 4 (pg. 111)}
\begin{itemize}
    \item T1: rotation around z-axis by 45 degrees counterclockwise
    \item T2: rotation around y-axis by 45 degrees clockwise
    \item T3: reflects the graph over the x-axis
    \item T4: reflects the graph over the y-axis
    \item T5: turns x-coordinate into negative z-coordinate, turns z-coordinate into negative x-coordinate, and keeps the y-coordinate the same
\end{itemize}

\section{Chapter 6 Exercise 5 (pg. 111)}
\begin{flalign*}
    P &= \left[ \begin{array}{cc}
        1 & 0 \\
        0 & 0
    \end{array} \right] \\
    P^2 &= \left[ \begin{array}{cc}
        1 & 0 \\
        0 & 0
    \end{array} \right] \left[ \begin{array}{cc}
        1 & 0 \\
        0 & 0
    \end{array} \right] = \left[ \begin{array}{cc}
        1\cdot 1 + 0 \cdot 0 & 1\cdot 0 + 0 \cdot 0 \\
        0\cdot 1 + 0 \cdot 0 & 0\cdot0 + 0\cdot 0
    \end{array} \right] = \left[ \begin{array}{cc}
        1 & 0 \\
        0 & 0
    \end{array} \right] = P \\
    Pe &= \lambda e \\
    \det(P - \lambda I) &= \det(\left[ \begin{array}{cc}
        1-\lambda & 0 \\
        0 & -\lambda
    \end{array} \right]) = (1-\lambda)(-\lambda) - 0 = 0 \Rightarrow \lambda \in \{ 1 \} \\
    \left[ \begin{array}{cc}
        1 & 0 \\
        0 & 0
    \end{array} \right] \left[ \begin{array}{cc}
        x \\ y
    \end{array} \right] &= 1 \cdot \left[ \begin{array}{cc}
        x \\ y
    \end{array} \right] \\
    e &= \left\{ \alpha \left[ \begin{array}{cc}
        x \\ 0
    \end{array} \right] : \alpha \in \mathbb{R} \backslash \{0\} \right\}
\end{flalign*}

\section{Chapter 6 Exercise 6 (pg. 111)}
\begin{flalign*}
    P &= \left[ \begin{array}{cc}
        1 & 0 \\
        0 & 1
    \end{array} \right] \\
    P^2 &= \left[ \begin{array}{cc}
        1 & 0 \\
        0 & 1
    \end{array} \right] \left[ \begin{array}{cc}
        1 & 0 \\
        0 & 1
    \end{array} \right] = \left[ \begin{array}{cc}
        1\cdot 1 + 0 \cdot 0 & 1\cdot 0 + 0 \cdot 1 \\
        0\cdot 1 + 1 \cdot 0 & 0\cdot0 + 1 \cdot 1
    \end{array} \right] = \left[ \begin{array}{cc}
        1 & 0 \\
        0 & 1
    \end{array} \right] = P \\
    Pe &= \lambda e \\
    \det(P - \lambda I) &= \det(\left[ \begin{array}{cc}
        1-\lambda & 0 \\
        0 & 1-\lambda
    \end{array} \right]) = (1-\lambda)(1-\lambda) - 0 = 0 \Rightarrow \lambda \in \{ 1 \} \\
    \left[ \begin{array}{cc}
        1 & 0 \\
        0 & 1
    \end{array} \right] \left[ \begin{array}{cc}
        x \\ y
    \end{array} \right] &= 1 \cdot \left[ \begin{array}{cc}
        x \\ y
    \end{array} \right] \\
    e &= \left\{ \alpha \left[ \begin{array}{cc}
        x \\ y
    \end{array} \right] : \alpha \in \mathbb{R} \backslash \{0\} \right\}
\end{flalign*}

\section{Chapter 6 Exercise 7 (pg. 111)}
\begin{flalign*}
    P &= \left[ \begin{array}{c}
        \cos(\theta) \\ \sin(\theta)
    \end{array}\right] \left[ \begin{array}{cc}
        \cos(\theta) & \sin(\theta)
    \end{array}\right] = \left[ \begin{array}{cc}
        \cos^2(\theta) & \cos(\theta)\sin(\theta) \\
        \cos(\theta)\sin(\theta) & \sin^2(\theta)
    \end{array}\right] \\
    P^2 &= \left[ \begin{array}{cc}
        \cos^2(\theta) & \cos(\theta)\sin(\theta) \\
        \cos(\theta)\sin(\theta) & \sin^2(\theta)
    \end{array}\right] \left[ \begin{array}{cc}
        \cos^2(\theta) & \cos(\theta)\sin(\theta) \\
        \cos(\theta)\sin(\theta) & \sin^2(\theta)
    \end{array}\right] \\
    &= \left[ \begin{array}{cc}
        \cos^4(\theta) + \cos^2(\theta)\sin^2(\theta) & \cos^3(\theta)\sin(\theta) + \cos(\theta)\sin^3(\theta) \\
        \cos^3(\theta)\sin(\theta) + \cos(\theta)\sin^3(\theta) & \cos^2(\theta)\sin^2(\theta) + \sin^4(\theta)
    \end{array}\right] \\
    &= \left[ \begin{array}{cc}
        \cos^2(\theta) (\cos^2(\theta) + \sin^2(\theta)) & \cos(\theta)\sin(\theta)(\cos^2(\theta) + \sin^2(\theta)) \\
        \cos(\theta)\sin(\theta)(\cos^2(\theta) + \sin^2(\theta)) & \sin^2(\theta)(\cos^2(\theta) + \sin^2(\theta))
    \end{array}\right] \\
    &= \left[ \begin{array}{cc}
        \cos^2(\theta) & \cos(\theta)\sin(\theta) \\
        \cos(\theta)\sin(\theta) & \sin^2(\theta)
    \end{array}\right] = P \\
    Pe &= \lambda e \\
    \det(P-\lambda I) &= \det(\left[ \begin{array}{cc}
        \cos^2(\theta) - \lambda & \cos(\theta)\sin(\theta) \\
        \cos(\theta)\sin(\theta) & \sin^2(\theta) - \lambda
    \end{array}\right]) = (\cos^2(\theta) - \lambda)(\sin^2(\theta) - \lambda) - \cos^2(\theta)\sin^2(\theta) \\
    &= \lambda^2 - \lambda(\cos^2(\theta) + \sin^2(\theta) = \lambda^2 - \lambda \Rightarrow \lambda \in \{ 0, 1 \} \\
    \left[ \begin{array}{cc}
        \cos^2(\theta) & \cos(\theta)\sin(\theta) \\
        \cos(\theta)\sin(\theta) & \sin^2(\theta)
    \end{array}\right] \left[ \begin{array}{cc}
        x \\ y
    \end{array} \right] &= 1 \cdot \left[ \begin{array}{cc}
        x \\ y
    \end{array} \right] \\
    e &= \left\{ \alpha \left[ \begin{array}{c}
    \end{array} \right] \right\}
\end{flalign*}

\section{Chapter 6 Exercise 8 (pg. 111)}
\begin{flalign*}
    \left[ \begin{array}{ccc|c}
        0 & 0 & 3 & 9 \\
        1 & 5 & -2 & 2 \\
        \frac{1}{3} & 2 & 0 & 3
    \end{array} \right] \\
    r_1 \leftarrow \frac{1}{3} r_1 \\
    \left[ \begin{array}{ccc|c}
        0 & 0 & 1 & 3 \\
        1 & 5 & -2 & 2 \\
        \frac{1}{3} & 2 & 0 & 3
    \end{array} \right] \\
    r_3 \leftarrow 3r_3 \\
    \left[ \begin{array}{ccc|c}
        0 & 0 & 1 & 3 \\
        1 & 5 & -2 & 2 \\
        1 & 6 & 0 & 9
    \end{array} \right] \\
    r_1 \text{ switch } r_3 \\
    \left[ \begin{array}{ccc|c}
        1 & 6 & 0 & 9 \\
        1 & 5 & -2 & 2 \\
        0 & 0 & 1 & 3
    \end{array} \right] \\
    r_2 \leftarrow r_2 + 2r_3 \\
    \left[ \begin{array}{ccc|c}
        1 & 6 & 0 & 9 \\
        1 & 5 & 0 & 8 \\
        0 & 0 & 1 & 3
    \end{array} \right] \\
    r_2 \leftarrow r_2 -r_1 \\
    \left[ \begin{array}{ccc|c}
        1 & 6 & 0 & 9 \\
        0 & -1 & 0 & -1 \\
        0 & 0 & 1 & 3
    \end{array} \right] \\
    r_2 \leftarrow -r_2 \\
    \left[ \begin{array}{ccc|c}
        1 & 6 & 0 & 9 \\
        0 & 1 & 0 & 1 \\
        0 & 0 & 1 & 3
    \end{array} \right] \\
    r_1 \leftarrow r_1 - 6r_2 \\
    \left[ \begin{array}{ccc|c}
        1 & 0 & 0 & 3 \\
        0 & 1 & 0 & 1 \\
        0 & 0 & 1 & 3
    \end{array} \right]
\end{flalign*}
\begin{itemize}
    \item Therefore, $x_1 = 3, x_2 = 1, x_3 = 3$
\end{itemize}

\section{Chapter }

\end{document}